\section{Conclusion}

This paper studied five-finger intent decoding from post-stroke HD-sEMG using a shared signal-processing and tensor-adaptation pipeline. Direct LSTM, CNN, and GNN baselines showed complementary strengths, with the LSTM leading exact-match accuracy and the GNN leading macro F1 and AUPRC. Subsequent CNN architecture search substantially improved the convolutional branch: CNN-Large reached the strongest offline aggregate performance, while CNN-Micro was selected for hardware deployment because it preserved most of that performance with 158K parameters and an estimated 154 KB INT8 footprint.

Healthy-to-impaired transfer learning further improved the CNN branch but produced mixed LSTM behavior, highlighting the architecture-specific nature of transfer in paretic sEMG decoding. Relative to prior work, the project is best understood as a hardware-aware post-stroke finger-intent study rather than a universal sEMG leaderboard result: it does not claim to surpass all healthy-subject or reduced-gesture benchmarks, but it does show that compact CNNs can approach or exceed internal ResNet teachers on a harder impaired-arm multilabel task. More broadly, the study shows that post-stroke finger-intent decoding benefits from treating signal processing, model family, transfer strategy, and deployment constraints as a single experimental design problem.
